\section{Involution classes in rank one}
\label{sec:rankone}

Here the dichotomy gives an explicit classification for an infinite family.
Throughout this section $q=p^f\ge5$ is odd. In $\PGL_2(q)$ the determinant modulo
squares defines a character with kernel $\PSL_2(q)$. Call an involution
\emph{interior} or \emph{exterior} according as its determinant is a square or a
non-square. These are the two involution classes of $\PGL_2(q)$.
The subgroup classification used below is Dickson's, in the form recorded by
Giudici~\cite[Theorems 2.2 and 3.5]{Giudici2007}; our conclusions concern the
commutator-tree levels for these specified leaf sets.

\begin{lemma}[Classes and subfields]
\label{lem:rankone-inputs}
The interior involutions generate $\PSL_2(q)$, and the exterior involutions
generate $\PGL_2(q)$. Every $A_5$ subgroup of $\PGL_2(q)$ is interior.
For a proper subfield $\F_{q_0}\subset\F_q$, exterior involutions of
$\PGL_2(q_0)$ remain exterior precisely when $[\F_q:\F_{q_0}]$ is odd.
\end{lemma}
\begin{proof}
The involutions of the simple group $\PSL_2(q)$ form a non-trivial normal
subset, so generate it. The subgroup generated by the exterior class is
normal in $\PGL_2(q)$. Its intersection with the simple socle is either the
socle or trivial; in the latter case its commutator with the socle is trivial,
contradicting the trivial centralizer of the socle in $\PGL_2(q)$. Hence it
contains the socle and an exterior element, and is the whole group.
The determinant character is trivial on the perfect group $A_5$.
Finally, if $m=[\F_q:\F_{q_0}]$, the parity of
$(q-1)/(q_0-1)=1+q_0+\cdots+q_0^{m-1}$ is the parity of $m$.
The square-character restriction therefore has the stated parity.
\end{proof}

\begin{theorem}[Rank-one classification]
\label{thm:rankone-classification}
For odd $q\ge5$ the following hold, with levels in the configurable model.
\begin{enumerate}[nosep,leftmargin=1.6em]
\item In $G=\PGL_2(q)$, the set of \emph{all} involutions has a
permanent obstruction.
\item In $G=\PGL_2(q)$, the exterior class has a permanent obstruction if and
only if there is a proper subfield $\F_{q_0}$ with $q_0\ge5$ and
$[\F_q:\F_{q_0}]$ odd. Otherwise its exact level is
\[
D=\begin{cases}3,&q\equiv\pm3\pmod8,\\2,&q\equiv\pm1\pmod8.\end{cases}
\]
\item In $G=\PSL_2(q)$, the involution class has a permanent obstruction if
and only if $G$ contains a \emph{proper} subgroup isomorphic to $A_5$, or
$\F_q$ has a proper subfield of size $q_0\ge5$. Otherwise its exact level is
\[
D=\begin{cases}3,&q\equiv\pm1\pmod8,\\2,&q\equiv\pm3\pmod8.\end{cases}
\]
In particular $\PSL_2(5)=A_5$ has level $2$. The $A_5$ obstruction in this
part occurs exactly when $q>5$ and either $p=5$ or $q\equiv\pm1\pmod5$.
\end{enumerate}
\end{theorem}
\begin{proof}
For (1), the interior leaves generate the proper non-solvable subgroup
$\PSL_2(q)$, so \Cref{thm:dichotomy} applies.

For (2), a proper exterior-generated subgroup lies in a maximal subgroup
other than $\PSL_2(q)$. Giudici's Theorem~3.5 lists these maximals: an affine
group $\F_q^+\rtimes\F_q^\times$, torus normalizers (dihedral groups), $S_4$
when $q=p\equiv\pm3\pmod8$, and $\PGL_2(q_0)$ when $q=q_0^r$ with $r$ an odd
prime. The affine group is metabelian, including over extension fields;
its involution-generated subgroups need not be ordinary dihedral groups.
Torus normalizers are also metabelian. The only non-solvable possibility
outside the socle is therefore a proper subfield group with $q_0\ge5$ and
odd degree. Such a group is indeed exterior-generated by
\Cref{lem:rankone-inputs}. Any odd proper degree has an odd prime divisor,
so the existence criterion agrees with this maximal-subgroup list.

In its absence, the possible subfield maximal with $q_0=3$ is
$\PGL_2(3)=S_4$, occurring for $q=3^r$, $r$ an odd prime. Together with the
prime-field $S_4$ case, these give an exterior-generated $S_4$ exactly in
the unobstructed cases $q\equiv\pm3\pmod8$.
(Indeed, this congruence forces $f$ odd; without a proper subfield of size
at least $5$ and odd degree, $f=1$ when $p\ge5$, and $f$ is an odd prime
when $p=3$.) Its determinant character is
the sign character, so its transpositions are exterior and its derived
length $3$ is attained. In the remaining cases all relevant maximals are
metabelian. The upper bound $2$ is attained: for $q\ge7$, take a split
torus generator $r:x\mapsto ax$ and an exterior involution $t:x\mapsto b/x$
($-b$ non-square). The two exterior involutions $t,r^2t$ generate
$\langle r^2,t\rangle$, a proper non-abelian dihedral group whose rotation
subgroup has order $(q-1)/2\ge3$. The case $q=5$ has the exterior $S_4$ and
was already covered. Apply \Cref{thm:dichotomy} for exactness.

For (3), Dickson's list says that any proper non-solvable subgroup lies in
an $A_5$ or a non-solvable subfield subgroup. A proper $A_5$ is generated
by its involutions. Every proper subfield of size $q_0\ge5$ supplies a
proper involution-generated $\PSL_2(q_0)$, irrespective of the degree;
even degree can also supply $\PGL_2(q_0)$. Conversely, the absence of these
two obstructions leaves affine and dihedral subgroups of derived length
at most $2$, and $A_4$ or $S_4$. The involutions in $A_4$ generate only
$V_4$. An $S_4$ occurs precisely for $q\equiv\pm1\pmod8$ and is generated
by involutions, giving level $3$ when unobstructed. Otherwise a
non-abelian dihedral torus normalizer gives level $2$: at least one of
$(q-1)/2,(q+1)/2$ is at least $3$, and its dihedral normalizer is proper.
This argument uses existence, not maximality, of these dihedral groups
at small $q$. At $q=5$ the $A_5$ is the ambient group and is not an
obstruction, in agreement with \Cref{cor:A5}. The stated $A_5$ embedding
condition is the corresponding part of Dickson's classification.
\end{proof}

\begin{corollary}[Prime fields]
\label{cor:rankone-prime}
For every prime $p\ge5$, the exterior involution class of $\PGL_2(p)$ mixes
at exact level $3$ for $p\equiv\pm3\pmod8$ and at exact level $2$ for
$p\equiv\pm1\pmod8$.
\end{corollary}

Thus $q=5,7,11$ give exterior levels $3,2,3$. At $q=11$ the involutions
of the $\PSL_2(11)$ ambient are instead obstructed by a proper $A_5$.
The extension degree matters: $q=25$ has exterior level $2$, whereas
$q=125$ contains the exterior-generated non-solvable $\PGL_2(5)$ and is
obstructed. \Cref{sec:quantitative-families} gives mixing rates as $q$ grows.

\begin{proposition}[An exact level beyond rank one]
\label{prop:psl33-level}
For $G=\PSL_3(3)$ and its involution class, the exact level is $5$.
\end{proposition}
\begin{proof}
Here $G=\SL_3(3)$ has trivial centre. Its involutions have eigenvalues
$1,-1,-1$, form a single conjugacy class, and generate $G$ by simplicity.
The maximal-subgroup list in the ATLAS~\cite{ATLASL33} consists of
two parabolic classes $3^2{:}\GL_2(3)$, one class of $S_4$, and $13{:}3$.

Write a point parabolic as $P=V\rtimes L$, where
$V\cong\F_3^2$ and $L\cong\GL_2(3)$. In the Levi factor,
\[
 L'=\SL_2(3),\qquad \SL_2(3)'=Q_8,\qquad
 Q_8'=\langle z\rangle,\qquad z=-I.
\]
The element $z$ acts as negation on $V$, so $[V,\langle z\rangle]=V$.
Consequently the derived series of $P$ is
\[
 P>V\rtimes\SL_2(3)>V\rtimes Q_8>
 V\rtimes\langle z\rangle>V>1,
\]
and $\operatorname{dl}(P)=5$.

The group $L$ is generated by
$s=\operatorname{diag}(-1,1)$ and the elementary transvections
$t=I+E_{12},I+E_{21}$. Since $sts=t^{-1}$ for either transvection,
both $s$ and $st$ are involutions, so $L$ is involution-generated.
For every $v\in V$, the product $(v z v^{-1})z$ is the translation
$2v$. Thus the involutions of $P$ also generate $V$, and hence $P$.
The line parabolic has the same properties by inverse transpose.

The remaining maximals are $S_4$, of derived length $3$, and
the odd-order group $13{:}3$. Every proper involution-generated
subgroup therefore has derived length at most $5$, and a parabolic
attains $5$.
Apply \Cref{thm:dichotomy}.
\end{proof}
